\FloatBarrier
\section{本章小结}
\label{sec:eigen-spectral-decomposition-summary}

本章开头问：给定方阵，在什么样的基下它最简单？答案取决于对基提出什么要求。三种要求依次加强，能达到的形式也依次退让。

只要求换基矩阵可逆时，目标是对角形。被作用只伸缩的方向是特征向量，特征值是特征多项式$p_{\bm A}(z)=\det(z\bm I-\bm A)$的根，与基的选取无关；Cayley--Hamilton定理说明$p_{\bm A}(\bm A)=\bm0$，所以$\bm A$的高次幂都能用前$n$个幂表示。

$n$阶矩阵实可对角化，当且仅当它有$n$个线性无关的实特征向量；等价地，特征多项式分解为实一次因子，且每个特征值的几何重数等于代数重数。这时$\bm A^k=\bm V\bm\Lambda^k\bm V^{-1}$，反复作用化为特征值的乘方。失败有两种方式：实特征值不够，如平面旋转；特征向量不够，如$\left[\begin{smallmatrix}0.8&1\\0&0.8\end{smallmatrix}\right]$，其幂含有随$k$增长的多项式因子。

要求换基矩阵正交时，特征基还必须两两垂直。能这样对角化的实矩阵恰好是对称矩阵，谱分解$\bm S=\bm Q\bm\Lambda\bm Q^\top$把作用写成沿互相垂直方向的伸缩。重特征值对应的特征向量可以在特征空间内任意旋转，但特征空间和谱投影由$\bm S$唯一确定，$\bm S=\sum_j\mu_j\bm P_j$；二次型在特征坐标下没有交叉项，所以定号性由特征值的符号决定。

对一般方阵，正交换基只能退到分块上三角的实Schur形：一阶对角块对应实特征值，二阶对角块对应一对共轭复特征值及其实不变平面。Schur形能取成分块对角的恰好是正规矩阵；这时每一步的最大放大倍数等于谱半径$\rho(\bm A)$。非正规矩阵的特征方向不正交，即使可以对角化，一步的放大倍数也可能远大于谱半径。

因此，特征方向不足时，应退到Schur形这类三角化表示；特征方向足够但不正交时，可逆换基仍能对角化，但长度和夹角在新坐标中失真，范数与谱半径也可能分离。第~\ref{chap:condition-and-spectral-perturbation}章将看到，这种不正交还使特征值对扰动更敏感；第~\ref{chap:linear-recurrence-stability}章将看到，它可能使反复作用先放大、再衰减。

本章的所有分解都把空间映到自身，并让输入和输出共用一组基。矩形矩阵和非正规方阵都需要另一种表示：第~\ref{chap:singular-value-decomposition}章让输入、输出各取一组标准正交基，构造对任何矩阵都存在的奇异值分解。
